\documentclass[11pt]{article}
\usepackage[T1]{fontenc}
\usepackage{textcomp}
\usepackage[margin=0.75in]{geometry}
\usepackage{amsmath,amssymb,amsthm}
\usepackage{array,booktabs}
\usepackage{hyperref}
\setlength{\parindent}{0pt}
\newtheorem{theorem}{Theorem}
\theoremstyle{definition}
\newtheorem{definition}{Definition}
\title{Flow Proof Artifact: prop-24-complements-similar-to-whole}
\author{Generated by \texttt{flow doc proof}}
\date{Flow Proof Book}
\begin{document}
\maketitle
In any parallelogram the complements of the parallelograms about the diameter are similar both to the whole and to one another.\\[0.5em]
\textbf{Source.} euclid — Elements, Book VI, Proposition 24\\[0.5em]
\bigskip
\noindent{\large \textbf{Derived fact 1.} \textit{In any parallelogram the complements of the parallelograms about the diameter are similar both to the whole and to one another} \hfill the Euclidean plane \cdot Euclid Book VI \cdot Proposition 24: parallelograms about a diameter are similar to the whole and to one another}\par
\smallskip\noindent\textit{Coordinate: the Euclidean plane \cdot Euclid Book VI \cdot Proposition 24: parallelograms about a diameter are similar to the whole and to one another}\par
\smallskip\noindent\textit{Source: euclid — Elements, Book VI, Proposition 24}\par
\smallskip\noindent\textit{Built on: proposition 34: complements of parallelograms about the diameter are equal, for Book I of the Elements in on the Euclidean plane, proposition 4: equiangular triangles have proportional corresponding sides, for Euclid Book VI on the Euclidean plane}\par
\medskip
\noindent\textbf{Goal.} In any parallelogram the complements of the parallelograms about the diameter are similar both to the whole and to one another\par
\noindent$\displaystyle \text{complements similar to whole}$\par
\medskip
\noindent
\begin{tabular}{@{} >{\raggedright\arraybackslash}p{0.06\textwidth} >{\raggedright\arraybackslash}p{0.47\textwidth} >{\raggedleft\arraybackslash}p{0.41\textwidth} @{}}
\textbf{\#} & \textbf{Proof} & \textbf{Mathematics} \\
\midrule
\textbf{1.} & We prove that proposition 24: parallelograms about a diameter are similar to the whole and to one another for Euclid Book VI the Euclidean plane. & \\[0.45em]
\textbf{2.} & Let ABCD = parallelogram. & \\[0.45em]
\textbf{3.} & Let BD = diameter. & \\[0.45em]
\textbf{4.} & Let comp1 = complement\_1. & \\[0.45em]
\textbf{5.} & From step 2, step 3, and step 4, we can deduce that comp1 \textasciitilde{} ABCD. & $\displaystyle comp1 ~ ABCD$ \\[0.45em]
\textbf{6.} & We invoke the derived fact governing Book I of the Elements in the Euclidean plane: proposition 34: complements of parallelograms about the diameter are equal, for Book I of the Elements in on the Euclidean plane. & \\[0.45em]
\textbf{7.} & We invoke the derived fact governing Euclid Book VI the Euclidean plane: proposition 4: equiangular triangles have proportional corresponding sides, for Euclid Book VI on the Euclidean plane. & \\[0.45em]
\textbf{8.} & From step 6 and step 7, this implies complements similar to whole. Hence proven. & $\displaystyle \text{complements similar to whole}$ \\[0.45em]
\bottomrule
\end{tabular}
\label{thm:Geometry-Euclid Book VI-proposition_24_parallelograms_about_a_diameter_are_similar_to_the_whole_and_to_one_another}
\par\medskip\hrule
\end{document}
